% Figure from MATH 493 notes, section 10. Compile with the notes preamble.
\begin{tikzpicture}[
    vx/.style={circle, draw=figB, thick, fill=white, minimum size=15pt, inner sep=0pt, font=\small},
    ax/.style={figR, thin, dashed},
    lb/.style={font=\scriptsize, figR}]
  \def\R{1.5}
  \fill[figB!12] (90:\R) -- (210:\R) -- (330:\R) -- cycle;
  \draw[figB, thick] (90:\R) -- (210:\R) -- (330:\R) -- cycle;
  % reflection axes: the axis through vertex i fixes i and swaps the other two
  \draw[ax] (90:2.0) -- (270:1.05);
  \draw[ax] (210:2.0) -- (30:1.05);
  \draw[ax] (330:2.0) -- (150:1.05);
  \node[lb, anchor=south] at (90:2.0) {$(2\,3)$};
  \node[lb, anchor=north east, inner sep=1pt] at (210:2.0) {$(1\,3)$};
  \node[lb, anchor=north west, inner sep=1pt] at (330:2.0) {$(1\,2)$};
  % rotation by 120 degrees counterclockwise: 1 -> 2 -> 3 -> 1
  \draw[figB, thick, -{Stealth[length=5pt]}] (118:1.9) arc (118:182:1.9);
  \node[font=\scriptsize, figB, anchor=south east, inner sep=1pt] at (150:2.0) {$(1\,2\,3)$};
  \node[vx] at (90:\R)  {$1$};
  \node[vx] at (210:\R) {$2$};
  \node[vx] at (330:\R) {$3$};
\end{tikzpicture}
